\section{The quartic straight-line response}
\label{sec:quartic}

The first nonlocal response about the straight line is quartic in the sources (Proposition~\ref{prop:quadratic}
has no cubic term). We analyse its structure, validate the strong-coupling input, and state what it can and
cannot provide.

\subsection{Operator content of the tilt channel}

Deform the internal direction to $n(\sigma)=(j_a(\sigma),\sqrt{1-j^2})$ on $S^5$, $a=1,\dots,5$, at $|\dot x|=1$. The
coupling is $\exp\int[j_a\Phi_a+(\sqrt{1-j^2}-1)\Phi_6]$, and odd orders vanish by $SO(5)$. The $O(j^4)$ term of
$\log Z$ has four pieces:
\begin{equation}
\frac1{4!}\!\int^4\!jjjj\,\dvev{\Phi_a\Phi_b\Phi_c\Phi_d}_{\rm conn}
-\frac14\!\int^3\!j_aj_bj^2\dvev{\Phi_a\Phi_b\Phi_6}
+\frac18\!\int^2\!j^2j^2\dvev{\Phi_6\Phi_6}_{\rm conn}
-\frac18\!\int\!j^4\dvev{\Phi_6}.
\end{equation}
$\Phi_6$ inserted on the half-BPS line is the lowest non-protected operator \cite{PolchinskiSully,BGT}, with
\begin{equation}
\Delta_6=1+\frac{\lambda}{4\pi^2}+O(\lambda^2),\qquad
\Delta_6=2-\frac5{\sqrt\lambda}+\frac{295}{24\lambda}-\frac{305}{16\lambda^{3/2}}+\Big(\frac{351845}{13824}-\frac{75}2\zeta(3)\Big)\frac1{\lambda^2}+\dots
\end{equation}
at weak and strong coupling respectively \cite{BGT,GRT,FerreroMeneghelli,GGJ}, and at strong coupling
$\Phi_6\leftrightarrow y^ay^a$. Remarks: $\vev{\Phi_6}$ vanishes in a conformal scheme but is a linear divergence in a
cutoff scheme, needed to make the tree-level quartic term finite; the $\Phi_6$ terms are contact-type (two source
times coincide), and at weak coupling $\dvev{\Phi_6\Phi_6}$ is the leading quartic term, $O(\lambda)$, giving
$j_a(t)\int\dd t'\,j(t')^2|t-t'|^{-2\Delta_6}$; the connected four-point function is $O(\lambda^2)$ at weak and
$O(1/\sqrt\lambda)$ at strong coupling in unit normalization, to be multiplied by $(2B)^2$.

\subsection{Structural propositions}

\begin{proposition}[the generalized free field is linear]
\label{prop:gff}
In the $\lambda\to\infty$ limit the defect fields are generalized free, $[\Phi(t),\Phi(0)]$ is a c-number, and every
nested commutator vanishes. The cubic retarded response at separated source times is therefore entirely the connected
four-point function, $O(1/\sqrt\lambda)$: equivalently, $\log Z=-S_{\rm cl}[n]+O(\lambda^0)$ and its quartic term is
the tree Witten diagram of the sphere fluctuations on AdS$_2$.
\end{proposition}

\begin{proposition}[memory]
\label{prop:memory}
The connected function at $O(1/\sqrt\lambda)$ contains $\log|\chi|$ and $\log|1-\chi|$ with nonzero coefficients
($-2/5$, $-1$, $+1$ at $\chi\to1$ in the singlet, symmetric and antisymmetric channels), whereas the generalized-free
part is rational. The Lorentzian continuation of a rational function gives contact-supported commutators; that of
the logarithms gives step-function discontinuities. The cubic response therefore has memory at $O(1/\sqrt\lambda)$,
and since the $\log\chi$ terms are $\gamma\log\chi$ from $\chi^{\Delta_0+\gamma/\sqrt\lambda}$, that memory is tied to
the anomalous dimensions, which resum to power laws with coupling-dependent exponents.
\end{proposition}

\begin{proposition}[no scale]
\label{prop:noscale}
Every retarded response kernel about the straight line is homogeneous in the time differences (the cubic tilt kernel
has degree $-4$); anomalous-dimension logarithms appear as $\log(t_{ij}/t_{kl})$; a constant background tilt is an
$SO(6)$ rotation of the half-BPS line. Hence no fixed delay, and in particular no hierarchy at $\log 2,\log3,\dots$,
can arise about the straight line at any order. A time-dependent background, a curved contour, a temperature or a
smeared observable would introduce a scale.
\end{proposition}

The Lorentzian cubic kernel itself, which needs all Wightman orderings (the continuation across $\chi=0,1,\infty$ with
$\ii\varepsilon$ prescriptions), was not computed.

\subsection{Validation of the strong-coupling input}

The tree-level connected four-point functions of the $\Delta=1$ fields $y^a$ (the superprimary of the displacement
multiplet) and of the $\Delta=2$ fields $x^i$ were taken from \cite{GRT} through an automated fetch summary and were
tested rather than trusted. With $\chi=t_{12}t_{34}/(t_{13}t_{24})$, $G=G^{(0)}+G^{(1)}/\sqrt\lambda$ and the $S,T,A$
structures, crossing $1\leftrightarrow2$ requires $G_{S,T}(\chi)=G_{S,T}(\chi/(\chi-1))$, $G_A(\chi)=-G_A(\chi/(\chi-1))$,
and crossing $1\leftrightarrow3$ requires $G(\chi)=(\chi/(1-\chi))^{2\Delta}R\,G(1-\chi)$ with
$R_5=\big[\begin{smallmatrix}1/5&28/25&-4/5\\1/2&3/10&1/2\\-1/2&7/10&1/2\end{smallmatrix}\big]$ for $SO(5)$ and
$R_3=\big[\begin{smallmatrix}1/3&10/9&-2/3\\1/2&1/6&1/2\\-1/2&5/6&1/2\end{smallmatrix}\big]$ for $SO(3)$, both with $R^2=1$.

\emph{$\Delta=1$.} The transcribed functions satisfy both crossings at both orders to $10^{-10}$;
$G_S^{(1)}=-2\chi^2\log\chi-\frac{43}{30}\chi^2+O(\chi^3\log\chi)$ with no constant and no linear term, so the identity is
not corrected; and the $\chi^n\log\chi$ coefficients divided by the generalized-free OPE coefficients $2/5$, $3/5$,
$-1$ give $\gamma_S=-5$ (the strong-coupling $\Delta_6$), $\gamma_T(\Delta_0=4)=-5$ and $\gamma_A(\Delta_0=3)=-4$, with no
$\chi^2\log\chi$ in the symmetric channel as the protected $\Delta=2$ operator requires. The transcription is validated.
Peeling off blocks order by order, a referee context found $\gamma_S(n)=-(2n^2+3n+5)$, $\gamma_T(n)=-(2n^2+3n)$,
$\gamma_A(n)=-(2n^2+5n+4)$, to be compared with \cite{GRT}.

\emph{$\Delta=2$.} The generalized-free part passes; the transcribed tree functions violate 16 of 18 crossing relations
and give $\gamma_T(\Delta_0=4)=-40/9$ instead of $-5$. That transcription is unreliable and is not used. The
displacement four-point function is in any case fixed by the superconformal Ward identities from the $\Delta=1$
function \cite{LMM,FerreroMeneghelli}.

For all-equal internal indices $G_{1111}=G_S+\tfrac85G_T$ and $G^{(1)}_{1111}=-2\chi^2\log\chi+\dots$, so the quartic
tilt integrand $G^{(1)}/(t_{12}^2t_{34}^2)$ is integrable at coincident points.

\subsection{Verdict}

The quartic response is the first nonlocal one on the straight line and its memory exponent, set by
$\Delta_6(\lambda)\in(1,2)$, is coupling dependent: a fixed protected dimension gives no variable exponent, but a
non-protected one does. However, the response is nonlinear, and the ordered tests presuppose a linear channel; a linear
channel would need a background supplying a scale; by Proposition~\ref{prop:noscale} no channel about the straight line
has fixed delays, so the prime-delay test cannot be passed in principle; and a dictionary $\lambda\leftrightarrow\omega$
would be an adjustment of the theory to the target. No arithmetic comparison is run.

\subsection*{Ledger}
\begin{center}\small
\begin{tabular}{@{}p{0.55\textwidth}p{0.40\textwidth}@{}}
\toprule
Statement & Status\\
\midrule
Operator content of the quartic tilt response & Proof (expansion); referee confirmed\\
$\Delta_6(\lambda)$ expansions & Literature, transcribed\\
Proposition~\ref{prop:gff} & Proof\\
Proposition~\ref{prop:memory} & Proof sketch (discontinuity structure); kernel not computed\\
Proposition~\ref{prop:noscale} & Proof (dilatations)\\
Crossing matrices & Proof; referee confirmed by index sums\\
$\Delta=1$ transcription validated & Numerical verification (93 cases)\\
$\Delta=2$ transcription unreliable & Numerical falsification\\
Anomalous-dimension patterns & Referee observation; compare with \cite{GRT}\\
\bottomrule
\end{tabular}
\end{center}
